\section{}

\paragraph{1099年}\eventanchor{XIA-1099-REGNAL-YEAR}{XIA-1099-REGNAL-YEAR}　西夏以李乾顺在位第13年作为诏令、赋役与外交文书的纪年框架。账本所列《宋史》卷485〈夏国传〉；《辽史》《金史》相关本纪；黑水城西夏文文书与考古报告为本条来源起点；该条可固定编年位置，具体执行、规模和地方社会影响仍须保留来源差异。

\paragraph{1100年}\eventanchor{XIA-1100-REGNAL-YEAR}{XIA-1100-REGNAL-YEAR}　西夏以李乾顺在位第14年作为诏令、赋役与外交文书的纪年框架。账本所列《宋史》卷485〈夏国传〉；《辽史》《金史》相关本纪；黑水城西夏文文书与考古报告为本条来源起点；该条可固定编年位置，具体执行、规模和地方社会影响仍须保留来源差异。

\paragraph{1101年}\eventanchor{XIA-1101-REGNAL-YEAR}{XIA-1101-REGNAL-YEAR}　西夏以李乾顺在位第15年作为诏令、赋役与外交文书的纪年框架。账本所列《宋史》卷485〈夏国传〉；《辽史》《金史》相关本纪；黑水城西夏文文书与考古报告为本条来源起点；该条可固定编年位置，具体执行、规模和地方社会影响仍须保留来源差异。

\paragraph{1102年}\eventanchor{XIA-1102-REGNAL-YEAR}{XIA-1102-REGNAL-YEAR}　西夏以李乾顺在位第16年作为诏令、赋役与外交文书的纪年框架。账本所列《宋史》卷485〈夏国传〉；《辽史》《金史》相关本纪；黑水城西夏文文书与考古报告为本条来源起点；该条可固定编年位置，具体执行、规模和地方社会影响仍须保留来源差异。

\paragraph{1103年}\eventanchor{XIA-1103-REGNAL-YEAR}{XIA-1103-REGNAL-YEAR}　西夏以李乾顺在位第17年作为诏令、赋役与外交文书的纪年框架。账本所列《宋史》卷485〈夏国传〉；《辽史》《金史》相关本纪；黑水城西夏文文书与考古报告为本条来源起点；该条可固定编年位置，具体执行、规模和地方社会影响仍须保留来源差异。

\paragraph{1104年}\eventanchor{XIA-1104-REGNAL-YEAR}{XIA-1104-REGNAL-YEAR}　西夏以李乾顺在位第18年作为诏令、赋役与外交文书的纪年框架。账本所列《宋史》卷485〈夏国传〉；《辽史》《金史》相关本纪；黑水城西夏文文书与考古报告为本条来源起点；该条可固定编年位置，具体执行、规模和地方社会影响仍须保留来源差异。

\paragraph{1105年}\eventanchor{XIA-1105-REGNAL-YEAR}{XIA-1105-REGNAL-YEAR}　西夏以李乾顺在位第19年作为诏令、赋役与外交文书的纪年框架。账本所列《宋史》卷485〈夏国传〉；《辽史》《金史》相关本纪；黑水城西夏文文书与考古报告为本条来源起点；该条可固定编年位置，具体执行、规模和地方社会影响仍须保留来源差异。

\paragraph{1106年}\eventanchor{XIA-1106-REGNAL-YEAR}{XIA-1106-REGNAL-YEAR}　西夏以李乾顺在位第20年作为诏令、赋役与外交文书的纪年框架。账本所列《宋史》卷485〈夏国传〉；《辽史》《金史》相关本纪；黑水城西夏文文书与考古报告为本条来源起点；该条可固定编年位置，具体执行、规模和地方社会影响仍须保留来源差异。

\paragraph{1107年}\eventanchor{XIA-1107-REGNAL-YEAR}{XIA-1107-REGNAL-YEAR}　西夏以李乾顺在位第21年作为诏令、赋役与外交文书的纪年框架。账本所列《宋史》卷485〈夏国传〉；《辽史》《金史》相关本纪；黑水城西夏文文书与考古报告为本条来源起点；该条可固定编年位置，具体执行、规模和地方社会影响仍须保留来源差异。

\paragraph{1108年}\eventanchor{XIA-1108-REGNAL-YEAR}{XIA-1108-REGNAL-YEAR}　西夏以李乾顺在位第22年作为诏令、赋役与外交文书的纪年框架。账本所列《宋史》卷485〈夏国传〉；《辽史》《金史》相关本纪；黑水城西夏文文书与考古报告为本条来源起点；该条可固定编年位置，具体执行、规模和地方社会影响仍须保留来源差异。

\paragraph{1109年}\eventanchor{XIA-1109-REGNAL-YEAR}{XIA-1109-REGNAL-YEAR}　西夏以李乾顺在位第23年作为诏令、赋役与外交文书的纪年框架。账本所列《宋史》卷485〈夏国传〉；《辽史》《金史》相关本纪；黑水城西夏文文书与考古报告为本条来源起点；该条可固定编年位置，具体执行、规模和地方社会影响仍须保留来源差异。

\paragraph{1110年}\eventanchor{XIA-1110-REGNAL-YEAR}{XIA-1110-REGNAL-YEAR}　西夏以李乾顺在位第24年作为诏令、赋役与外交文书的纪年框架。账本所列《宋史》卷485〈夏国传〉；《辽史》《金史》相关本纪；黑水城西夏文文书与考古报告为本条来源起点；该条可固定编年位置，具体执行、规模和地方社会影响仍须保留来源差异。

\paragraph{1111年}\eventanchor{XIA-1111-REGNAL-YEAR}{XIA-1111-REGNAL-YEAR}　西夏以李乾顺在位第25年作为诏令、赋役与外交文书的纪年框架。账本所列《宋史》卷485〈夏国传〉；《辽史》《金史》相关本纪；黑水城西夏文文书与考古报告为本条来源起点；该条可固定编年位置，具体执行、规模和地方社会影响仍须保留来源差异。

\paragraph{1112年}\eventanchor{XIA-1112-REGNAL-YEAR}{XIA-1112-REGNAL-YEAR}　西夏以李乾顺在位第26年作为诏令、赋役与外交文书的纪年框架。账本所列《宋史》卷485〈夏国传〉；《辽史》《金史》相关本纪；黑水城西夏文文书与考古报告为本条来源起点；该条可固定编年位置，具体执行、规模和地方社会影响仍须保留来源差异。

\paragraph{1113年}\eventanchor{XIA-1113-REGNAL-YEAR}{XIA-1113-REGNAL-YEAR}　西夏以李乾顺在位第27年作为诏令、赋役与外交文书的纪年框架。账本所列《宋史》卷485〈夏国传〉；《辽史》《金史》相关本纪；黑水城西夏文文书与考古报告为本条来源起点；该条可固定编年位置，具体执行、规模和地方社会影响仍须保留来源差异。

\paragraph{1114年}\eventanchor{XIA-1114-REGNAL-YEAR}{XIA-1114-REGNAL-YEAR}　西夏以李乾顺在位第28年作为诏令、赋役与外交文书的纪年框架。账本所列《宋史》卷485〈夏国传〉；《辽史》《金史》相关本纪；黑水城西夏文文书与考古报告为本条来源起点；该条可固定编年位置，具体执行、规模和地方社会影响仍须保留来源差异。

\paragraph{1115年}\eventanchor{XIA-1115-EVENT-067}{XIA-1115-EVENT-067}　金建国改变西夏北方外交环境。账本所列《宋史》卷485〈夏国传〉；《辽史》《金史》相关本纪；黑水城西夏文文书与考古报告为本条来源起点；该条可固定编年位置，具体执行、规模和地方社会影响仍须保留来源差异。

\paragraph{1115年}\eventanchor{XIA-1115-REGNAL-YEAR}{XIA-1115-REGNAL-YEAR}　西夏以李乾顺在位第29年作为诏令、赋役与外交文书的纪年框架。账本所列《宋史》卷485〈夏国传〉；《辽史》《金史》相关本纪；黑水城西夏文文书与考古报告为本条来源起点；该条可固定编年位置，具体执行、规模和地方社会影响仍须保留来源差异。

\paragraph{1116年}\eventanchor{XIA-1116-REGNAL-YEAR}{XIA-1116-REGNAL-YEAR}　西夏以李乾顺在位第30年作为诏令、赋役与外交文书的纪年框架。账本所列《宋史》卷485〈夏国传〉；《辽史》《金史》相关本纪；黑水城西夏文文书与考古报告为本条来源起点；该条可固定编年位置，具体执行、规模和地方社会影响仍须保留来源差异。

\paragraph{1117年}\eventanchor{XIA-1117-REGNAL-YEAR}{XIA-1117-REGNAL-YEAR}　西夏以李乾顺在位第31年作为诏令、赋役与外交文书的纪年框架。账本所列《宋史》卷485〈夏国传〉；《辽史》《金史》相关本纪；黑水城西夏文文书与考古报告为本条来源起点；该条可固定编年位置，具体执行、规模和地方社会影响仍须保留来源差异。

\paragraph{1118年}\eventanchor{XIA-1118-REGNAL-YEAR}{XIA-1118-REGNAL-YEAR}　西夏以李乾顺在位第32年作为诏令、赋役与外交文书的纪年框架。账本所列《宋史》卷485〈夏国传〉；《辽史》《金史》相关本纪；黑水城西夏文文书与考古报告为本条来源起点；该条可固定编年位置，具体执行、规模和地方社会影响仍须保留来源差异。

\paragraph{1119年}\eventanchor{XIA-1119-REGNAL-YEAR}{XIA-1119-REGNAL-YEAR}　西夏以李乾顺在位第33年作为诏令、赋役与外交文书的纪年框架。账本所列《宋史》卷485〈夏国传〉；《辽史》《金史》相关本纪；黑水城西夏文文书与考古报告为本条来源起点；该条可固定编年位置，具体执行、规模和地方社会影响仍须保留来源差异。

\paragraph{1120年}\eventanchor{XIA-1120-REGNAL-YEAR}{XIA-1120-REGNAL-YEAR}　西夏以李乾顺在位第34年作为诏令、赋役与外交文书的纪年框架。账本所列《宋史》卷485〈夏国传〉；《辽史》《金史》相关本纪；黑水城西夏文文书与考古报告为本条来源起点；该条可固定编年位置，具体执行、规模和地方社会影响仍须保留来源差异。

\paragraph{1121年}\eventanchor{XIA-1121-REGNAL-YEAR}{XIA-1121-REGNAL-YEAR}　西夏以李乾顺在位第35年作为诏令、赋役与外交文书的纪年框架。账本所列《宋史》卷485〈夏国传〉；《辽史》《金史》相关本纪；黑水城西夏文文书与考古报告为本条来源起点；该条可固定编年位置，具体执行、规模和地方社会影响仍须保留来源差异。

\paragraph{1122年}\eventanchor{XIA-1122-REGNAL-YEAR}{XIA-1122-REGNAL-YEAR}　西夏以李乾顺在位第36年作为诏令、赋役与外交文书的纪年框架。账本所列《宋史》卷485〈夏国传〉；《辽史》《金史》相关本纪；黑水城西夏文文书与考古报告为本条来源起点；该条可固定编年位置，具体执行、规模和地方社会影响仍须保留来源差异。

\paragraph{1123年}\eventanchor{XIA-1123-EVENT-068}{XIA-1123-EVENT-068}　金在灭辽过程中与西夏建立新的外交关系。账本所列《宋史》卷485〈夏国传〉；《辽史》《金史》相关本纪；黑水城西夏文文书与考古报告为本条来源起点；该条可固定编年位置，具体执行、规模和地方社会影响仍须保留来源差异。

\paragraph{1123年}\eventanchor{XIA-1123-REGNAL-YEAR}{XIA-1123-REGNAL-YEAR}　西夏以李乾顺在位第37年作为诏令、赋役与外交文书的纪年框架。账本所列《宋史》卷485〈夏国传〉；《辽史》《金史》相关本纪；黑水城西夏文文书与考古报告为本条来源起点；该条可固定编年位置，具体执行、规模和地方社会影响仍须保留来源差异。

\paragraph{1124年}\eventanchor{XIA-1124-REGNAL-YEAR}{XIA-1124-REGNAL-YEAR}　西夏以李乾顺在位第38年作为诏令、赋役与外交文书的纪年框架。账本所列《宋史》卷485〈夏国传〉；《辽史》《金史》相关本纪；黑水城西夏文文书与考古报告为本条来源起点；该条可固定编年位置，具体执行、规模和地方社会影响仍须保留来源差异。

\paragraph{1125年}\eventanchor{XIA-1125-REGNAL-YEAR}{XIA-1125-REGNAL-YEAR}　西夏以李乾顺在位第39年作为诏令、赋役与外交文书的纪年框架。账本所列《宋史》卷485〈夏国传〉；《辽史》《金史》相关本纪；黑水城西夏文文书与考古报告为本条来源起点；该条可固定编年位置，具体执行、规模和地方社会影响仍须保留来源差异。

\paragraph{1126年}\eventanchor{XIA-1126-REGNAL-YEAR}{XIA-1126-REGNAL-YEAR}　西夏以李乾顺在位第40年作为诏令、赋役与外交文书的纪年框架。账本所列《宋史》卷485〈夏国传〉；《辽史》《金史》相关本纪；黑水城西夏文文书与考古报告为本条来源起点；该条可固定编年位置，具体执行、规模和地方社会影响仍须保留来源差异。

\paragraph{1127年}\eventanchor{XIA-1127-REGNAL-YEAR}{XIA-1127-REGNAL-YEAR}　西夏以李乾顺在位第41年作为诏令、赋役与外交文书的纪年框架。账本所列《宋史》卷485〈夏国传〉；《辽史》《金史》相关本纪；黑水城西夏文文书与考古报告为本条来源起点；该条可固定编年位置，具体执行、规模和地方社会影响仍须保留来源差异。

\paragraph{1128年}\eventanchor{XIA-1128-REGNAL-YEAR}{XIA-1128-REGNAL-YEAR}　西夏以李乾顺在位第42年作为诏令、赋役与外交文书的纪年框架。账本所列《宋史》卷485〈夏国传〉；《辽史》《金史》相关本纪；黑水城西夏文文书与考古报告为本条来源起点；该条可固定编年位置，具体执行、规模和地方社会影响仍须保留来源差异。

\paragraph{1129年}\eventanchor{XIA-1129-REGNAL-YEAR}{XIA-1129-REGNAL-YEAR}　西夏以李乾顺在位第43年作为诏令、赋役与外交文书的纪年框架。账本所列《宋史》卷485〈夏国传〉；《辽史》《金史》相关本纪；黑水城西夏文文书与考古报告为本条来源起点；该条可固定编年位置，具体执行、规模和地方社会影响仍须保留来源差异。

\paragraph{1130年}\eventanchor{XIA-1130-REGNAL-YEAR}{XIA-1130-REGNAL-YEAR}　西夏以李乾顺在位第44年作为诏令、赋役与外交文书的纪年框架。账本所列《宋史》卷485〈夏国传〉；《辽史》《金史》相关本纪；黑水城西夏文文书与考古报告为本条来源起点；该条可固定编年位置，具体执行、规模和地方社会影响仍须保留来源差异。

\paragraph{1131年}\eventanchor{XIA-1131-REGNAL-YEAR}{XIA-1131-REGNAL-YEAR}　西夏以李乾顺在位第45年作为诏令、赋役与外交文书的纪年框架。账本所列《宋史》卷485〈夏国传〉；《辽史》《金史》相关本纪；黑水城西夏文文书与考古报告为本条来源起点；该条可固定编年位置，具体执行、规模和地方社会影响仍须保留来源差异。

\paragraph{1132年}\eventanchor{XIA-1132-REGNAL-YEAR}{XIA-1132-REGNAL-YEAR}　西夏以李乾顺在位第46年作为诏令、赋役与外交文书的纪年框架。账本所列《宋史》卷485〈夏国传〉；《辽史》《金史》相关本纪；黑水城西夏文文书与考古报告为本条来源起点；该条可固定编年位置，具体执行、规模和地方社会影响仍须保留来源差异。

\paragraph{1133年}\eventanchor{XIA-1133-REGNAL-YEAR}{XIA-1133-REGNAL-YEAR}　西夏以李乾顺在位第47年作为诏令、赋役与外交文书的纪年框架。账本所列《宋史》卷485〈夏国传〉；《辽史》《金史》相关本纪；黑水城西夏文文书与考古报告为本条来源起点；该条可固定编年位置，具体执行、规模和地方社会影响仍须保留来源差异。

\paragraph{1134年}\eventanchor{XIA-1134-REGNAL-YEAR}{XIA-1134-REGNAL-YEAR}　西夏以李乾顺在位第48年作为诏令、赋役与外交文书的纪年框架。账本所列《宋史》卷485〈夏国传〉；《辽史》《金史》相关本纪；黑水城西夏文文书与考古报告为本条来源起点；该条可固定编年位置，具体执行、规模和地方社会影响仍须保留来源差异。

\paragraph{1135年}\eventanchor{XIA-1135-REGNAL-YEAR}{XIA-1135-REGNAL-YEAR}　西夏以李乾顺在位第49年作为诏令、赋役与外交文书的纪年框架。账本所列《宋史》卷485〈夏国传〉；《辽史》《金史》相关本纪；黑水城西夏文文书与考古报告为本条来源起点；该条可固定编年位置，具体执行、规模和地方社会影响仍须保留来源差异。

\paragraph{1136年}\eventanchor{XIA-1136-REGNAL-YEAR}{XIA-1136-REGNAL-YEAR}　西夏以李乾顺在位第50年作为诏令、赋役与外交文书的纪年框架。账本所列《宋史》卷485〈夏国传〉；《辽史》《金史》相关本纪；黑水城西夏文文书与考古报告为本条来源起点；该条可固定编年位置，具体执行、规模和地方社会影响仍须保留来源差异。

\paragraph{1137年}\eventanchor{XIA-1137-REGNAL-YEAR}{XIA-1137-REGNAL-YEAR}　西夏以李乾顺在位第51年作为诏令、赋役与外交文书的纪年框架。账本所列《宋史》卷485〈夏国传〉；《辽史》《金史》相关本纪；黑水城西夏文文书与考古报告为本条来源起点；该条可固定编年位置，具体执行、规模和地方社会影响仍须保留来源差异。

\paragraph{1138年}\eventanchor{XIA-1138-REGNAL-YEAR}{XIA-1138-REGNAL-YEAR}　西夏以李乾顺在位第52年作为诏令、赋役与外交文书的纪年框架。账本所列《宋史》卷485〈夏国传〉；《辽史》《金史》相关本纪；黑水城西夏文文书与考古报告为本条来源起点；该条可固定编年位置，具体执行、规模和地方社会影响仍须保留来源差异。

\paragraph{1139年}\eventanchor{XIA-1139-EVENT-069}{XIA-1139-EVENT-069}　李仁孝即位。账本所列《宋史》卷485〈夏国传〉；《辽史》《金史》相关本纪；黑水城西夏文文书与考古报告为本条来源起点；该条可固定编年位置，具体执行、规模和地方社会影响仍须保留来源差异。

\paragraph{1139年}\eventanchor{XIA-1139-REGNAL-YEAR}{XIA-1139-REGNAL-YEAR}　西夏以李乾顺在位第53年作为诏令、赋役与外交文书的纪年框架。账本所列《宋史》卷485〈夏国传〉；《辽史》《金史》相关本纪；黑水城西夏文文书与考古报告为本条来源起点；该条可固定编年位置，具体执行、规模和地方社会影响仍须保留来源差异。

\paragraph{1140年}\eventanchor{XIA-1140-REGNAL-YEAR}{XIA-1140-REGNAL-YEAR}　西夏以李仁孝在位第1年作为诏令、赋役与外交文书的纪年框架。账本所列《宋史》卷485〈夏国传〉；《辽史》《金史》相关本纪；黑水城西夏文文书与考古报告为本条来源起点；该条可固定编年位置，具体执行、规模和地方社会影响仍须保留来源差异。

\paragraph{1141年}\eventanchor{XIA-1141-REGNAL-YEAR}{XIA-1141-REGNAL-YEAR}　西夏以李仁孝在位第2年作为诏令、赋役与外交文书的纪年框架。账本所列《宋史》卷485〈夏国传〉；《辽史》《金史》相关本纪；黑水城西夏文文书与考古报告为本条来源起点；该条可固定编年位置，具体执行、规模和地方社会影响仍须保留来源差异。

\paragraph{1142年}\eventanchor{XIA-1142-REGNAL-YEAR}{XIA-1142-REGNAL-YEAR}　西夏以李仁孝在位第3年作为诏令、赋役与外交文书的纪年框架。账本所列《宋史》卷485〈夏国传〉；《辽史》《金史》相关本纪；黑水城西夏文文书与考古报告为本条来源起点；该条可固定编年位置，具体执行、规模和地方社会影响仍须保留来源差异。

\paragraph{1143年}\eventanchor{XIA-1143-REGNAL-YEAR}{XIA-1143-REGNAL-YEAR}　西夏以李仁孝在位第4年作为诏令、赋役与外交文书的纪年框架。账本所列《宋史》卷485〈夏国传〉；《辽史》《金史》相关本纪；黑水城西夏文文书与考古报告为本条来源起点；该条可固定编年位置，具体执行、规模和地方社会影响仍须保留来源差异。

\paragraph{1144年}\eventanchor{XIA-1144-EVENT-070}{XIA-1144-EVENT-070}　西夏与金以册封、贡赐和边境安排维持和平。账本所列《宋史》卷485〈夏国传〉；《辽史》《金史》相关本纪；黑水城西夏文文书与考古报告为本条来源起点；该条可固定编年位置，具体执行、规模和地方社会影响仍须保留来源差异。

\paragraph{1144年}\eventanchor{XIA-1144-REGNAL-YEAR}{XIA-1144-REGNAL-YEAR}　西夏以李仁孝在位第5年作为诏令、赋役与外交文书的纪年框架。账本所列《宋史》卷485〈夏国传〉；《辽史》《金史》相关本纪；黑水城西夏文文书与考古报告为本条来源起点；该条可固定编年位置，具体执行、规模和地方社会影响仍须保留来源差异。
